\section{Preliminary results and overview of previous research}

% This chapter describes the physical properties of superconducting qubits, which are fundamental to quantum information
% processing using superconducting circuits, and their control methods. First, we clarify how macroscopic quantum systems,
% specifically superconducting circuits, are quantized and behave as qubits through the derivation of the Hamiltonian.
% Next, we explain the mechanisms of control electronics (Z-control and XY-control) for manipulating single-qubit states.
% Furthermore, we detail the operational principles and adiabatic implementation of the CZ gate and
% conclude with an overview of relevant prior work to clarify the position of this research.
After we briefly overview the basics of quantum bits, this chapter introduces some fundamental theoretical frameworks and
reviews prior methodologies essential for understanding the optimization of control pulses in superconducting quantum
systems. The Derivative Removal by Adiabatic Gate (DRAG) formalism, Schrieffer-Wolff transformations, and Riccati
engineering serve as the baseline for the advanced waveform optimizations proposed in this study. In exploring
these control methodologies, we first consider the case of the one-qubit X gate before later extending our analysis to
the two-qubit CZ gate architecture.

\subsection{Introduction to the Quantum Bit}
The transition from classical to quantum information processing represents a
fundamental paradigm shift in computational theory. In classical computing, the foundational unit of information is
bit which exists in a mutually exclusive, deterministic state of either 0 or 1. Quantum
computing replaces this discrete entity with the quantum bit, or qubit, an entity governed by the principles of
quantum mechanics. Rather than being confined to a single state, a qubit can exist in a coherent superposition of
both states simultaneously, unlocking exponentially larger state spaces and enabling computational speedups for
specific classes of problems, such as integer factorization and molecular simulation.

Mathematically, the state of a single, isolated qubit is represented as a state vector in a two-dimensional
complex Hilbert space. Using Dirac notation, the computational basis states are denoted as $\vert{}0\rangle$ and
$\vert{}1\rangle$, corresponding to the classical 0 and 1. A pure qubit state $\vert{}\psi\rangle$ is expressed as a
linear combination of these basis states:$$\vert{}\psi\rangle = \alpha\vert{}0\rangle + \beta\vert{}1\rangle$$Here,
$\alpha$ and $\beta$ are complex probability amplitudes satisfying
$\vert{}\alpha\vert{}^2 + \vert{}\beta\vert{}^2 = 1$.
This constraint allows the pure state of a qubit to be mapped onto the surface of a unit sphere in three-dimensional
space, known as the Bloch sphere. On the Bloch sphere, the poles represent the basis states $\vert{}0\rangle$ and
$\vert{}1\rangle$, while points on the equator represent equal superpositions with varying relative phases. This
geometric representation is indispensable for visualizing single-qubit quantum gates as rotations about the sphere's axes.

As a quantum mechanical system, a qubit's dynamics are governed by the time-dependent Schrödinger equation.
Consequently, any control applied to the system is mathematically represented by the Hamiltonian operator driving this
evolution. In formula, \cite{sakurai2020modern}
$$
  \psi(t) = \mathcal{T}\left(\exp\left[\int_0^T\,iH(t)\,\dd t\right]\right)\psi(0).
$$
\subsection{Hamiltonian in the rotating frame for a qudit}
A transmon qubit is a widely used superconducting qubit. Because a transmon is inherently an anharmonic oscillator, it possesses higher excited states beyond
$\ket{0}$ and $\ket{1}$.
Consequently, it acts as a qudit---a quantum system with $d$ levels---rather than an ideal qubit, which
inevitably introduces leakage into unwanted noncomputational states.

To model this leakage analytically, we review the derivation of the qudit Hamiltonian truncated at the first $\Nmax+1$
levels of the transmon by following Ref.~\cite{99999}.
In passing, we derive the Hamiltonian without the Rotating Wave Approximation (RWA)\cite{Fujii2017}.

The Hamiltonian of a transmon qudit in the lab frame is \cite{wong2025hardware}
\begin{equation}
  \Hlab = 4E_c(\num - \num_g)^2 - E_J \cos(\hat{\varphi}), \label{eq:lab}
\end{equation}
where $\hat{n}_g$ denotes the offset charge, $\hat{n}$ the number operator, $E_C$ the charging energy, $E_J$ the Josephson
energy, and $\varphi$ the phase operator.
From the (non-degenerate) perturbation theory, we find that the eigenstates of this Hamiltonian are
harmonics with non-uniform energy spacings.
We denote the $k$-th excited state as $\ket{k}$ and its $\varphi$-representation $\psi_k(\varphi)$.
Now we consider the expansion of $\num$ in $\ket{k}\bra{m}\, (k\geq 0,\, m\geq 0)$.
Note that $k$'s parity coincides with $\psi_k$'s.
Since $\num$ flips the sign, the coefficient
\begin{equation}
  \lambda_{k,m}=\bra{k}\num\ket{m}\,=\,\int_{-\pi}^\pi\,\psi_k^*(\varphi)\left(-i\right)
\end{equation}
is zero unless $k$ and $m$ are of opposite parity.
Hence $\hat{n}$ can be expressed as
\begin{eqnarray}\nonumber
  \hat{n}&=&\sum_{k=0}^{N_{{\rm{max}}}-1}\bigg[\lambda_{k,k+1}\ketbra{k}{k+1}\\
    &+&\sum_{j=1} \lambda_{k,k+2j+1}\ketbra{k}{k+2j+1}\bigg]+\rm{h.c}.
\end{eqnarray}

If we apply the change-of-basis $V$ on the Hilbert space, the Hamiltonian changes as
\begin{equation}
  H^\prime = VHV^\dagger + i\dot{V}V^\dagger,
  \label{eq:framechange}
\end{equation}
where the last term is called a gauge term.
Applying the diagonal frame transformation with
\begin{equation}
  R\updag = \exp\left(i\wdrive t\sumkk[\Nmax] k\ketbra{k}{k}\right)
\end{equation}
yields the rotating frame Hamiltonian
\begin{align}
  \Hrot= & \sum_{k=0}^{N_{{\rm{max}}}}{\Delta}_{k}\ket{k}\bra{k}+ \OmegaRF
  \Biggr[\nonumber                                                                                         \\
         & \sum_{k=0}^{N_{\rm{max}}-1}\sum_{j=1}\lambda_{k,k+2j+1}\ketbra{k}{k+2j+1}e^{-(2j+1)i\omega_d t}
  \nonumber                                                                                                \\
         & +
    \sum_{k=0}^{N_{{\rm{max}}}-1}\lambda_{k,k+1}\ketbra{k}{k+1}e^{-i\omega_d t}
    +\rm{h.c.}
    \Biggr]
  ,
  \label{eq:rotatingframe}
\end{align}
where ${\Delta}_k = \omega_k - k\omega_d$, $\omega_d$ is the drive frequency and $\omega_k = E_k - E_0$.
Henceforth, if we simply write $\Delta$, this means $\Delta_2$.

If we apply RWA entirely, \eref{eq:rotatingframe} reduces to
\begin{eqnarray}\nonumber
  \Hrwa &=& \sum_{j=1}^{\Nmax} \left(\Delta_j+j\delta(t)\right)\ket{j}\bra{j}\\
  &+&\sum_{j=1}^{\Nmax}\lambda_j\bigg[\frac{\Omega_x(t)}{2} \hat{\sigma}^x_{j-1, j}+\frac{\Omega_y(t)}{2} \hat{\sigma}^y_{j-1, j}\bigg],
  \label{eq:Hrwa}
\end{eqnarray}
where $\lambda_j = \lambda_{j-1,j}$, $\hat{\sigma}^x_{j-1,j} = \ket{j-1}\bra{j} + \ket{j}\bra{j-1}$, and
$\hat{\sigma}^y_{j-1,j} = -i\ket{j-1}\bra{j} + i\ket{j}\bra{j-1}$.

The matricial form when $\Nmax=3$ is
\begin{equation}
  \footnotesize
  \left(\begin{array}{cccc}
      0                                                & \frac{1}{2} \lambda_1 (\Omegax(t)-i \Omegay(t)) & 0                                               & 0                                               \\
      \frac{1}{2} \lambda_1  (\Omegax(t)+i \Omegay(t)) & \delta (t)+\Delta_1                             & \frac{1}{2} \lambda_2 (\Omegax(t)-i \Omegay(t)) & 0                                               \\
      0                                                & \frac{1}{2} \lambda_2 (\Omegax(t)+i \Omegay(t)) & 2 \delta (t)+\Delta_2                           & \frac{1}{2} \lambda_3 (\Omegax(t)-i \Omegay(t)) \\
      0                                                & 0                                               & \frac{1}{2} \lambda_3 (\Omegax(t)+i \Omegay(t)) & 3 \delta (t)+\Delta_3                           \\
    \end{array}\right).
  \normalsize
  \label{eq:rwa_mat}
\end{equation}
Including the first order corrections of RWA, the Hamiltonian becomes
\begin{align}
  \Hrwa[1] = & \sum_{k=0}^{N_{{\rm{max}}}}\tilde{\Delta}_{k}\ket{k}\bra{k}+               %\OmegaRF
  \frac{1}{2}\left(\Omega_x-i\Omega_y\right)
  \Biggr[\nonumber                                                                        \\
             & \sum_{k=0}^{N_{\rm{max}}-1}\lambda_{k,k+3}\ketbra{k}{k+3}e^{-2i\omega_d t}
  \nonumber                                                                               \\
             & +
    \sum_{k=0}^{\Nmax-1}\lambda_{k,k+1}\ketbra{k}{k+1}
    \Biggr]
  +\rm{h.c.}
  % .
  \label{eq:rwa_first_order}
\end{align}
where, for brevity, the last \rm{h.c.} only refers off-diagonal components.

In this study, we consider up to the second excited state in our analysis.
% where $\Delta_{k}=\omega_{k}-k\omega_{d}$ is the detuning between the $k$-th level with the $k$-th driving frequency
% harmonic. In this frame, all coupling terms, except for $\ket{k}\leftrightarrow\ket{k+1}$, and all the counter-rotating
% components oscillate rapidly. Therefore, we can neglect those small contributions, leading to the Hamiltonian in
% \eref{eq:rwa_first_order}.

\subsection{Approximation of $\lambda$}
In the QHO approximation, $\hat{n}$ is expressed in terms of the annihilation and creation operators:
$$\hat{n} \approx in_{\text{zpf}}(\hat{a}^\dagger - \hat{a})$$
Recall $\hat{a} = \sum_{k=0}^\infty\,\sqrt{k+1}\ket{k+1}\bra{k}$.
$$
  \bra{k+1}\hat{a}\ket{k} = \sqrt{k+1}\quad
  \bra{k+1}\hat{a}^\dagger\ket{k} = 0
$$
Hence,
$$
  \lambda_{k,k+1} \approx \bra{k+1}\hat{n}\ket{k} = in_{\text{zpf}}\sqrt{k+1}
$$
The prefactor $in_{\text{zpf}}$  could be regarded as absorbed by the drive amplitude $\OmegaRF$.
Hence we simply set
$$
  \lambda_{k,k+1} \approx \sqrt{k+1}.
$$
For $\lambda_1$ is entirely aborbed into $\OmegaRF$, it is exactly (not approximately) $\lambda_1=1$.

In simulations, we set the values of $\lambda$ accordiingly.
\subsection{Schrieffer-Wolff transformation}
\label{subsec:swt}
In \eref{eq:framechange}, $V$ is assumed to be a unitary operator.
we now consider a scenario where the transformation operator $V$ is parameterized by an  anti-Hermitian operator $S$.
Specifically, $V$ is defined exponentially as:
$$
  V = \exp(S),
$$
we can expand the Hamiltonian in a power series of $S$ as
\begin{equation}
  H^\prime = i\dot{S} + H + [S,H] + \frac{1}{2}[S,[S,H]] + \cdots. \quad(\text{Hadamard's lemma})
  \label{eq:swt}
\end{equation}
It is important to emphasize that this expansion is mathematically rigorous and holds generally true for any generator
$S$ and Hamiltonian $H$, provided the system is defined within a finite-dimensional Hilbert space.
While standard perturbative techniques often approximate the dynamics by truncating this series early, exploring a
non-perturbative, strong-driving regime requires a more exacting approach. In such strongly driven systems, premature
truncation introduces significant inaccuracies.
% In the DRAG scheme, this evaluation of higher order terms is called superlinear corrections.


\subsection{DRAG}
\label{subsec:drag}
DRAG is an extremely effective method to reduce leakage.
Gambetta and Motzoi derived the DRAG Y correction formula
\begin{equation}
  \Omega_y = -\frac{1}{\Delta}\dot{\Omega}_x
  \label{eq:ydrag}
\end{equation}
in \cite{Gambetta2011}.
% This formula even holds in non-perturbative settings.
% % Applying a frame change with $V=\exp(S)$, $S=-i\lambda_2\dfrac{\Omegax}{2\Delta_2}\hat{\sigma}^y_{1,2}$ to the Hamiltonian \ref{eq:Hrwa} ,
% % the new $\ket{1}\bra{2}$ component reads
% Consider a frame change $V$
% \begin{equation}
%   V = \exp\left[-i\lambda_2\dfrac{\Omegax}{2\Delta_2}\hat{\sigma}^y_{1,2}\right]
%   = \cos(\phi) -i\hat{\sigma}^y_{1,2}\sin(\phi)
%   \label{eq:first_frame_change}
% \end{equation}
% with $\phi = \lambda_2\dfrac{\Omegax}{2\Delta_2}$.
% Then the resulting $\ket{1}\bra{2}$ component is
% $$
%   -\sin(2\phi)\dfrac{\delta(t)+\Delta}{2}+\cos(2\phi)\dfrac{\lambda \Omegax}{2}-i\dfrac{\lambda}{2}\left(\dfrac{\Omega_x^\prime(t)}{\Delta}+\Omega_y(t)\right).
% $$
% Setting $\Omegay$ as in \eref{eq:ydrag} still reduces this component.
They have also showed that AC Stark shift can be compensated by phase ramping
\begin{equation}
  \delta(t) = \dfrac{4-\lambda^2}{4\Delta}\,\Omegax^2.
  \label{eq:zdrag}
\end{equation}
Collectively \eref{eq:ydrag} and \eref{eq:zdrag} are called DRAG conditions.
If these two are applied in weak driving regime, the Hamiltonian becomes
\begin{equation}
  \hat{\mathcal{H}}_1(t) \approx \frac{\Omega_x}{2}\hat{\sigma}_{0,1}^x + \frac{\lambda_2 \Omega_x^2}{8\Delta_2}\hat{\sigma}_{0,2}^x - \Omega_x^2 \left[ \frac{(2\Delta_2 - \Delta_3)}{8\Delta_2^2} \right] \lambda_2 \lambda_3 \hat{\sigma}_{1,3}^x
  \label{eq:Hdrag}
\end{equation}
with the following boundary conditions to satisfy
$$
  \Omega_x(t) = \dot{\Omega}_x(t)=0 \quad \text{at }t=0,\,T
$$
\subsubsection{DRAG-1}
Hamiltonian \eref{eq:Hdrag} is remained with two photon transitions between $\ket{0}\leftrightarrow\ket{2}$ and
$\ket{1}\leftrightarrow\ket{3}$.
SWTs to further reduce these transitions is applicable.
Now we introduce a recursive DRAG method called DRAG-1 \cite{Li2025} or R1D \cite{99999}:

Applying Schrieffer-Wolff transformations with $S_2 = -i\lambda_2\Omega_1^2/(8\Delta_2^2)\hat{\sigma}^y_{0,2}$ and
$S_3 = i\lambda_2\dot{\Omega}_1\Omega_1/(4\Delta_2^3)\hat{\sigma}^x_{0,2}$
to
\eref{eq:Hdrag} yields
$$
  \hat{\mathcal{H}}_\text{new}(t) = \frac{\Omega_x}{2}\hat{\sigma}_{0,1}^x + \frac{\lambda_2 \Omega_x^2}{8\Delta_2}\hat{\sigma}_{0,2}^x - \Omega_x^2 \left[ \frac{(2\Delta_2 - \Delta_3)}{8\Delta_2^2} \right] \lambda_2 \lambda_3 \hat{\sigma}_{1,3}^x
  - \lambda_2\left[\right]\sigma^x_{0,2}
$$
Therefore,
\begin{equation}
  \Omega_x(t) = \sqrt{\Omega_1^2 + \frac{2}{\Delta_2^2} (\dot{\Omega}_1^2 + \ddot{\Omega}_1 \Omega_1)}
  \label{eq:drag1}
\end{equation}
$\Omega_1$ has to satisfy the following boundary conditions
$$
  \Omega_1 = 0 \quad \text{at }t=0,\,T
$$

In case of $X$ gate, the $X$ gate condition
\begin{equation}
  \int_0^T\,\Omega_x(t)\,\dd t = \pi
  \label{eq:xgate_condition}
\end{equation}
gets a slight modification if we consider up to the third degree in the frame transformation \eref{eq:first_frame_change}
$$
  \Omega_x^\prime(t) = \Omega_x(t) - \dfrac{(4-\lambda_2^2)\Omega_x^3(t)}{8\Delta^2}.
$$
(Corrections of this sort in the DRAG scheme are called superlinear corrections.)


\subsection{NPAD}
\label{subsec:npad}
A lot of efforts for leakage reduction and fidelity improvements have been done in the perturbative regime (i.e.
$\Omega/\Delta \ll 1$).
One of the framework advocated for the non-perturbative strong driving regime is NPAD \cite{NPAD2022}.
NPAD is an iterative method to diagonalize a nonperturbative
Hamiltonian by applying static Givens rotations until the desired accuracy is achieved.
While this method is only effective if the Hamiltonian has a rather simple form and only applicable to time-independent
Hamiltonian, theoretically this is a generalization of recursive SWT in the perturbative regime.


\subsection{Riccati engineering}
Although widely known as a qunatum control method, Riccati engineering is not yet applied to time-dependent one qubit
gate controls.
Given a Hamiltonian
$$H(t) = \begin{pmatrix} H_{11}(t) & H_{12}(t) \\ H_{21}(t) & H_{22}(t) \end{pmatrix},$$
we consider block diagonalization of $H(t)$ into $2\times 2$ and $1\times 1$ partition according to the transformation
\eref{eq:framechange}.

The motivation is that if we could block diagonalize $H$ in this way, the dynamics starting with $\ket{0}$
in this new frame will cause zero leakage into the non-computational subspace.
As long as the off-diagonal terms are zero.
% \begin{boxA}
%   {\bf Probelm } \\
%   Given a $3\times 3$ matrix $H$, block diagonalize $H$ into $2\times 2$ and $1\times 1$
%   compartments. Namely, if we write
%   apply a frame change transformation \eref{eq:framechange} with a unitary matrix $U(t)$ and
%   make $H_{12}(t)$ and $H_{21}(t)$ $0$.
%   % $$H'(t) = U(t)H(t)U^\dagger(t) + i\dot{U}(t)U^\dagger(t)$.
% \end{boxA}
The solution to this problem is given in \cite{Gardas2010},
$$U(t) = \begin{pmatrix} (I_2 + X^\dagger X)^{-1/2} & X^\dagger (I_{n-2} + X X^\dagger)^{-1/2} \\ -X (I_2 + X^\dagger X)^{-1/2} & (I_{n-2} + X X^\dagger)^{-1/2} \end{pmatrix}$$
with $X(t)=(x_1(t)\:\:\: x_2(t))$ satisfying
\begin{equation}
  i\dot{X} = H_{21} + H_{22}X - XH_{11} - XH_{12}X \label{eq:riccati}
\end{equation}
Under this frame change, the resulting $H_{11}^\prime(t)$ is given by
\begin{equation}
  H^\prime_{11}(t) = \left(I_2+X^\dagger X\right)^{1/2} (H_{11}+H_{12}X)(I_2+X^\dagger X)^{-1/2}
  - (I_2+X^\dagger X)^{1/2} \partial_t\left[(I_2+X^\dagger X)^{-1/2}\right]
\end{equation}
The form of this Hamiltonian is not very easy to analyze.
The cases in which the Hamiltonian become time-dependent are investigated in \cite{Riccati2010,Gardas2010}.

\subsection {CZ gate}
From now, we start dealing with the two qubit CZ gate.
Note that we use different notations for X and CZ gates.
\subsubsection{Basics}
CZ gate is a two-qubit gate that applies a phase flip to the target qubit if the control qubit is in the state
$\ket{1}$. In the computational basis $\{\vert{}00\rangle, \vert{}01\rangle, \vert{}10\rangle, \vert{}11\rangle\}$, the
ideal unitary evolution is represented by a diagonal matrix:
\begin{equation}
  U_{\mathrm{CZ}}=\begin{pmatrix}
    1 & 0 & 0 & 0  \\
    0 & 1 & 0 & 0  \\
    0 & 0 & 1 & 0  \\
    0 & 0 & 0 & -1
  \end{pmatrix}
\end{equation}
\subsubsection{SQUID loop and Z-control}
In superconducting quantum architectures, the ability to dynamically tune the transition frequency of a qubit is
essential for realizing the adiabatic CZ gate.
This tunability is achieved by replacing the single Josephson junction of a standard
transmon with a superconducting quantum interference device (SQUID) loop, which is then threaded by a localized magnetic
flux controlled via a dedicated microwave transmission line known as the Z-control line (or flux line).
\subsubsection{Qubit Architecture}
Consider a system of two capacitively coupled transmon qubits, one fixed-frequency (QB2) and the other frequency tunable (QB1).
The simplified dispersive Jaynes-Cummings Hamiltonian for this system can be modeled as:
$$H = \sum_{i=1,2} \left( \omega_i a_i^\dagger a_i + \frac{\alpha_i}{2} a_i^\dagger a_i^\dagger a_i a_i \right) + g(a_1^\dagger + a_1)(a_2^\dagger + a_2) ,$$
where $g$ denotes the coupling strength between the two qubits, $a_i$ creation operator for qubit $i$, and $\alpha_i$
the anharmonicity of qubit $i$.
\subsubsection{Adiabatic Implementation}
% The adiabatic CZ gate is realized as syncing the QB1's frequency to QB2's, thereby creating an avoided crossing.
% By doing so, the energy of $\ket{11}$ and $\ket{20}$ become comparable.
To implement the gate, the frequency of Qubit 1 is swept along a
prescribed trajectory $\omega_1(t)$ toward the avoided crossing, held or continuously varied for a duration $\tau_d$,
and then returned to its idle parking frequency.
In this process, $\ket{11}$ obtains/loses more phases than the other computational states, due to its stronger coupling
to $\ket{02}$.

By following \cite{Ding2025,Martinis2014} we review the theory behind the adiabatic implementaion of CZ gate and the
derivation of leakage formula \eref{eq:leakage}.
\begin{theorem}[Adiabatic Theorem]
  If a quantum system is initially prepared in a non-degenerate energy eigenstate of a time-dependent Hamiltonian, and
  if that Hamiltonian is varied sufficiently slowly, the system will evolve continuously to remain in the corresponding
  instantaneous eigenstate of the evolving Hamiltonian up to an accumulated phase factor, given that there is a gap
  between the eigenvalue and the rest of the Hamiltonian's spectrum.
\end{theorem}
Consider a two level system with the following Hamiltonian
$$
  H(t)=\eps(t)\sigma_z + \Delta\sigma_x.
$$
According to the adiabatic theorem, if we prepare the state $\ket{0}$ and slowly vary the Hamiltonian and then retrace the
change back, the resulting state is $\ket{0}$ with phase accumulation.
The instantaneous eigenstates are
\begin{equation}
  \vert{}\psi_-\rangle = \begin{bmatrix} -\sin(\theta(t)/2) \\ \cos(\theta(t)/2) \end{bmatrix} \quad \text{and} \quad \vert{}\psi_+\rangle = \begin{bmatrix} \cos(\theta(t)/2) \\ \sin(\theta(t)/2) \end{bmatrix},
\end{equation}
with eigenenergies
\begin{equation}
  E_- = -\frac{1}{2}\sqrt{\varepsilon(t)^2 + \Delta^2} \quad \text{and} \quad E_+ = \frac{1}{2}\sqrt{\varepsilon(t)^2 + \Delta^2}
\end{equation}
where
\begin{equation}
  \theta(t) = \arctan\frac{\Delta}{\eps(t)}.
  \label{eq:theta}
\end{equation}
We denote the energy difference of these two instantaneous eigenstates as
\begin{equation}
  \omega(t) = \sqrt{\varepsilon(t)^2 + \Delta^2}
  \label{eq:omega}
\end{equation}

\subsubsection{Calculation of leakage error} \label{subsec:leakage}
Consider a state
$$
  \ket{\psi} = \alpha \ket{\psi_{-}}+\beta\ket{\psi_+} = \alpha_0\ket{0}+\beta_0\ket{1}.,
$$
with rotational change of coefficients
\begin{align*}
  \alpha & = \alpha_0\cos\dfrac{\theta}{2} + \beta_0\sin\dfrac{\theta}{2}  \\
  \beta  & =  \beta_0\cos\dfrac{\theta}{2} - \alpha_0\sin\dfrac{\theta}{2}
\end{align*}
Differentiating, we have
\begin{align*}
  \dot{\alpha} & = \dfrac{-i\omega \alpha+\beta \dot{\theta}}{2} \\
  \dot{\beta}  & = \dfrac{i\omega \beta-\alpha \dot{\theta}}{2}  \\
\end{align*}
Write $\alpha = \cos \Theta/2$ and $\beta=e^{i\phi}\sin\Theta/2$, they form the "angle":
$$
  \alpha^* \beta = \dfrac{\sin\Theta e^{i\phi}}{2}
$$
Calculating the time derivative of $\alpha^*\beta$ finds us
$$
  \frac{d}{dt}(\sin\Theta e^{i\phi}) = i\omega \sin\Theta e^{i\phi}  \pm  \dot{\theta}\cos\Theta
$$
Plugging $\phi=\phi_{\text{initial}} + \int_0^t\,\omega(t^\prime)\,\dd t^\prime$ and a bit of tinkering yields
\begin{equation}
  % \frac{d}{dt}(\sin\Theta e^{i\phi}) = i\omega \sin\Theta e^{i\phi}  \pm  \dot{\theta}\cos\Theta
  P_e \approx \dfrac{1}{4}\, \left\lvert\, \int \,\dfrac{\dd\theta}{\dd t} e^{-i\int\omega(t^\prime)\,\dd t^\prime}\right\rvert^2
  \label{eq:leakage}
\end{equation}
% If we assume that the time derivative of $\ket{\psi_-}$ and $\ket{\psi_+}$ are small enough, $\ket{\psi}$ evolves accordingly to \begin{align*}
%   \dot{\alpha} & = -iE_- \alpha, \\
%   \dot{\beta}  & = -iE_+\beta,
% \end{align*}
% in other words,
% \begin{align*}
%   \dot{\alpha} & = \dot{\alpha_0}\cos\dfrac{\theta}{2} + \dot{\beta_0}\sin\dfrac{\theta}{2},  \\
%   \dot{\beta}  & =  \dot{\beta_0}\cos\dfrac{\theta}{2} - \dot{\alpha_0}\sin\dfrac{\theta}{2},
% \end{align*}


\subsubsection{Phase accumulation}
According to the adiabatic theorem, if the sweep rate is sufficiently slow compared to the energy gap at the avoided crossing,
the system will remain in its instantaneous eigenstate. The dynamic frequency excursion imparts a continuous deformation
on the $\vert{}11\rangle$ state, causing it to hybridize temporarily with $\vert{}02\rangle$ and acquire a dynamical phase.
The total accumulated phase $\Phi_{11}$ on the $\vert{}11\rangle$ state is the time integral of its instantaneous eigenenergy
$E_{11}(t)$:$$\Phi_{11} = -\frac{1}{\hbar} \int_{0}^{\tau} E_{11}(t) dt$$
Simultaneously, single-qubit dynamical phases $\Phi_{01}$ and
$\Phi_{10}$ are acquired. The conditional phase $\phi_{\text{cond}}$ accumulated over the gate duration is calculated
by isolating the two-qubit interaction phase from the single-qubit contributions:$$\phi_{\text{cond}} = \Phi_{11} -
  \Phi_{01} - \Phi_{10} + \Phi_{00}$$The gate is successful when the pulse amplitude and duration are calibrated such
that $\phi_{\text{cond}} = \pi \pmod{2\pi}$.

\subsubsection{Nonlinear time transformation}
By applying a time transformation, we can simplify the leakage formula \eref{eq:leakage}.
Let the time variable of this new frame be $\tau$. Set
$$
  \Delta \dd \tau = \omega(t) \dd t.
$$
then
\begin{equation}
  \dd t = \dfrac{\Delta}{\omega(t)}\dd \tau = \sin \theta(t) \dd \tau.
  \label{eq:nonlinear_time_transform_differential}
\end{equation}
In integral form,
\begin{equation}
  t(\tau) = \int_0^\tau\,\sin\tilde{\theta}(\tau^\prime)\,\dd\tau^\prime
  \label{eq:nonliniear_time_transform}
\end{equation}
The virtue of this transformation is clear if we plug \eref{eq:nonlinear_time_transform_differential} into \eref{eq:leakage}
\begin{equation}
  P_e = \dfrac{1}{4}\, \left\lvert\, \int_0^\tau\,\dfrac{\dd\tilde{\theta}}{\dd\tau}\, e^{-i\Delta \tau}\,\dd \tau \right\rvert^2
  \label{eq:leakage_tau}
\end{equation}

\subsubsection{Slepian pulses}

In the context of pulse waveform optimization, the Slepian pulses, frequently referred to as discrete prolate spheroidal
sequences (DPSSs), represent a family of orthogonal pulses specifically formulated to solve the problem of maximal
energy concentration in both the time and frequency domains. Because Heisenberg's uncertainty principle dictates that a
pulse cannot be strictly bounded in both time and frequency simultaneously, Slepian, Landau, and Pollack developed these
sequences to answer a fundamental optimization question: how to optimally concentrate energy in one domain when the
pulse is strictly confined in the opposing domain.

To define this optimization rigorously, consider a discrete-time pulse of finite length, denoted as $x[n]$ for $n = 0,
  1, \dots, N-1$. This pulse is constrained to possess finite energy $E$, defined as:$$E = \sum_{n=0}^{N-1}
  \vert{}x[n]\vert{}^2 < \infty$$The frequency characteristics of the pulse are analyzed using its discrete-time Fourier
transform:$$X(e^{i\omega}) = \sum_{n=0}^{N-1} x[n]e^{-i\omega n}$$The objective is to maximize the fraction of the total
pulse energy that is contained within a specific frequency band $[-W, W]$, where $0 < W < \pi$. This fractional energy
concentration is represented by the ratio $\lambda \in [0, 1]$:$$\lambda = \frac{\int_{-W}^{W}
    \vert{}X(e^{i\omega})\vert{}^2 d\omega}{\int_{-\pi}^{\pi} \vert{}X(e^{i\omega})\vert{}^2 d\omega}$$Therefore, the core
optimization problem is to identify the specific sequence $x[n]$ of length $N$ that maximizes this eigenvalue
$\lambda$.3. The Eigenvalue SolutionThe sequences that solve this optimization problem, denoted as $v_n^{(k)}(N, W)$
where $k = 0, 1, \dots, N-1$ represents the pulse order, are the Slepian pulses. Here, $N$ is the pulse length and $W$
parameterizes the mainlobe width.These pulses are derived as the real solutions to the following system of
equations:$$\sum_{m=0}^{N-1} \frac{\sin 2\pi W(n - m)}{\pi(n - m)} v_m^{(k)}(N, W) = \lambda_k(N, W)v_n^{(k)}(N, W)$$For
the singularity condition where $n = m$, the fraction simplifies to $2W$.This system can be elegantly expressed as the
matrix eigenvalue problem $\mathbf{A}\mathbf{v}^{(k)} = \lambda_k \mathbf{v}^{(k)}$.The matrix elements are defined as
$A_{n,m} = \frac{\sin 2\pi W(n - m)}{\pi(n - m)}$.The eigenvector is $\mathbf{v}^{(k)} = [v_0^{(k)}(N, W), v_1^{(k)}(N,
  W), \dots, v_{N-1}^{(k)}(N, W)]^T$.4. Characteristics and Practical AdvantagesSolving the matrix equation yields $N$
distinct eigenvalues, which are conventionally ranked in descending order: $1 > \lambda_0 > \lambda_1 > \dots >
  \lambda_{N-1} > 0$.The sequence $v_n^{(0)}(N, W)$, associated with the strictly largest eigenvalue $\lambda_0$, is
designated as the first Slepian pulse.Each subsequent higher-order pulse maximizes the energy ratio $\lambda$ while
maintaining strict orthogonality to all preceding Slepian pulses.When evaluating their performance against simpler
shapes like rectangular or raised cosine pulses, Slepian pulses are distinguished by possessing both a relatively small
mainlobe width and low sidelobe amplitude.This dual efficiency makes them an exceptionally strong candidate for control
applications where a precise compromise between mainlobe width and sidelobe suppression is required.
